\chapter{Longest Decreasing Subsequence (LDS)}

\section{Definizione del problema}

\subsection{Istanza}
Sia data una sequenza $X = \langle x_1, x_2, \dots, x_n \rangle$ di $n$ numeri reali (o interi).

\subsection{Soluzione}
Una più lunga sottosequenza $Z = \langle x_{i_1}, x_{i_2}, \dots, x_{i_k} \rangle$ di $X$ (con $1 \le i_1 < i_2 < \dots < i_k \le n$) i cui elementi sono ordinati in senso strettamente decrescente:
\[
    x_{i_1} > x_{i_2} > \dots > x_{i_k}
\]

\section{Mancanza di informazione e definizione del coefficiente}
Approcciando il problema sui soli prefissi $X_i = \langle x_1, \dots, x_i \rangle$ senza ulteriori vincoli, si incorre in una \textbf{mancanza di informazione}: conoscendo unicamente la lunghezza della più lunga sottosequenza decrescente di $X_{i-1}$, non è noto quale sia il suo ultimo elemento e non è quindi possibile verificare se $x_i$ sia strettamente minore di esso.

Si definisce pertanto un problema ausiliario \textbf{vincolato} a terminare con l'elemento $x_i$.

\medskip
\noindent
\textbf{Definizione del coefficiente $c_i$:} Per ogni $i \in \{1, \dots, n\}$, $c_i$ è la lunghezza della più lunga sottosequenza strettamente decrescente del prefisso $X_i$ la quale è \textbf{vincolata a terminare con $x_i$}.

\section{Equazioni di ricorrenza}

\subsection{Caso base}
\[
    i = 1 \implies c_1 = 1
\]
\textbf{Giustificazione:} L'unica sottosequenza del prefisso $X_1 = \langle x_1 \rangle$ che termina con $x_1$ è $\langle x_1 \rangle$ stessa, la quale è banalmente decrescente e ha lunghezza $1$.

\subsection{Passo ricorsivo}
Per $i > 1$:
\[
    c_i = \max \{ c_j : 1 \le j < i \land x_j > x_i \} + 1 \quad (\text{con } \max \emptyset = 0)
\]
\textbf{Giustificazione:} Ogni sottosequenza strettamente decrescente di $X_i$ che termina con $x_i$ si ottiene accodando $x_i$ a una sottosequenza strettamente decrescente di un prefisso $X_j$ ($1 \le j < i$) terminante in $x_j$, purché sia rispettata la condizione di decrescenza $x_j > x_i$. Per massimizzare la lunghezza si sceglie il $j$ che massimizza $c_j$; se nessun elemento precedente è maggiore di $x_i$ (insieme vuoto), $x_i$ forma da solo una sottosequenza di lunghezza $0 + 1 = 1$.

\subsection{Soluzione del problema principale}
Poiché la più lunga sottosequenza decrescente dell'intera sequenza $X$ deve terminare con uno degli elementi $x_i$ ($i \in \{1, \dots, n\}$), la lunghezza ottima globale è:
\[
    \text{Soluzione} = \max_{1 \le i \le n} \{ c_i \}
\]

\section{Algoritmo Bottom-Up}
L'algoritmo utilizza un vettore $c[1 \dots n]$ per i coefficienti e un vettore dei predecessori $b[1 \dots n]$ per memorizzare l'indice $j$ che ha generato il massimo per $c[i]$ (oppure $0$ se l'insieme dei candidati è vuoto).

\begin{algorithm}[htbp]
\caption{Calcolo Bottom-Up per LDS}
\begin{algorithmic}[1]
\Procedure{LDS}{$X$}
    \State $n \gets |X|$
    \State $c \gets \text{nuovo array } c[1 \dots n]$
    \State $b \gets \text{nuovo array } b[1 \dots n]$
    \Comment{Caso base per $i = 1$}
    \State $c[1] \gets 1$
    \State $b[1] \gets 0$
    \State $m \gets c[1]$
    \State $\text{ind\_m} \gets 1$
    \Comment{Passo ricorsivo}
    \For{$i = 2 \textbf{ to } n$}
        \State $\text{tmp} \gets 0$
        \State $k \gets 0$
        \For{$j = 1 \textbf{ to } i - 1$}
            \If{$x_i < x_j \land c[j] > \text{tmp}$}
                \State $\text{tmp} \gets c[j]$
                \State $k \gets j$
            \EndIf
        \EndFor
        \State $c[i] \gets \text{tmp} + 1$
        \State $b[i] \gets k$
        \If{$c[i] > m$}
            \State $m \gets c[i]$
            \State $\text{ind\_m} \gets i$
        \EndIf
    \EndFor
    \State \Return $c, b, m, \text{ind\_m}$
\EndProcedure
\end{algorithmic}
\end{algorithm}

\section{Algoritmo di ricostruzione}
\begin{algorithm}[htbp]
\caption{Ricostruzione della sequenza LDS}
\begin{algorithmic}[1]
\Procedure{Print-LDS}{$i, b, X$}
    \If{$b[i] == 0$}
        \State \textbf{print} $x_i$
    \Else
        \State \Call{Print-LDS}{$b[i], b, X$}
        \State \textbf{print} $x_i$
    \EndIf
\EndProcedure
\end{algorithmic}
\end{algorithm}

\noindent
\textbf{Chiamata principale (Driver):} Una volta eseguito $\textsc{LDS}(X)$, si invoca $\Call{Print-LDS}{\text{ind\_m}, b, X}$.

\section{Analisi della complessità}
\begin{itemize}
    \item \textbf{Complessità spaziale:} I due vettori $c$ e $b$ hanno dimensione $n$, per cui lo spazio richiesto è $\Theta(n)$.
    \item \textbf{Complessità temporale (Bottom-Up):} Per ogni $i \in \{2, \dots, n\}$ il ciclo interno esegue $i - 1$ confronti di costo $\Theta(1)$:
    \[
        T(n) = \sum_{i=2}^{n} \Theta(i - 1) = \Theta\!\left(\frac{n(n-1)}{2}\right) = \Theta(n^2)
    \]
    \item \textbf{Complessità temporale (Ricostruzione):} La procedura risale la catena dei predecessori stampando esattamente $m \le n$ elementi in tempo $O(n)$.
\end{itemize}